% Options for packages loaded elsewhere
% Options for packages loaded elsewhere
\PassOptionsToPackage{unicode}{hyperref}
\PassOptionsToPackage{hyphens}{url}
%
\documentclass[
  ignorenonframetext,
  aspectratio=169,
  russian,
  aspectratio=169]{beamer}
\newif\ifbibliography
\usepackage{pgfpages}
\setbeamertemplate{caption}[numbered]
\setbeamertemplate{caption label separator}{: }
\setbeamercolor{caption name}{fg=normal text.fg}
\beamertemplatenavigationsymbolshorizontal
% remove section numbering
\setbeamertemplate{part page}{
  \centering
  \begin{beamercolorbox}[sep=16pt,center]{part title}
    \usebeamerfont{part title}\insertpart\par
  \end{beamercolorbox}
}
\setbeamertemplate{section page}{
  \centering
  \begin{beamercolorbox}[sep=12pt,center]{section title}
    \usebeamerfont{section title}\insertsection\par
  \end{beamercolorbox}
}
\setbeamertemplate{subsection page}{
  \centering
  \begin{beamercolorbox}[sep=8pt,center]{subsection title}
    \usebeamerfont{subsection title}\insertsubsection\par
  \end{beamercolorbox}
}
% Prevent slide breaks in the middle of a paragraph
\widowpenalties 1 10000
\raggedbottom
\AtBeginPart{
  \frame{\partpage}
}
\AtBeginSection{
  \ifbibliography
  \else
    \frame{\sectionpage}
  \fi
}
\AtBeginSubsection{
  \frame{\subsectionpage}
}
\usepackage{iftex}
\ifPDFTeX
  \usepackage[T1]{fontenc}
  \usepackage[utf8]{inputenc}
  \usepackage{textcomp} % provide euro and other symbols
\else % if luatex or xetex
  \usepackage{unicode-math} % this also loads fontspec
  \defaultfontfeatures{Scale=MatchLowercase}
  \defaultfontfeatures[\rmfamily]{Ligatures=TeX,Scale=1}
\fi
\usepackage{lmodern}

\usefonttheme{serif} % use mainfont rather than sansfont for slide text
\ifPDFTeX\else
  % xetex/luatex font selection
  \setmainfont[]{DejaVu Serif}
  \setsansfont[]{DejaVu Sans}
  \setmonofont[]{DejaVu Sans Mono}
\fi
% Use upquote if available, for straight quotes in verbatim environments
\IfFileExists{upquote.sty}{\usepackage{upquote}}{}
\IfFileExists{microtype.sty}{% use microtype if available
  \usepackage[]{microtype}
  \UseMicrotypeSet[protrusion]{basicmath} % disable protrusion for tt fonts
}{}


\usepackage{longtable,booktabs,array}
\newcounter{none} % for unnumbered tables
\usepackage{calc} % for calculating minipage widths
\usepackage{caption}
% Make caption package work with longtable
\makeatletter
\def\fnum@table{\tablename~\thetable}
\makeatother
\usepackage{graphicx}
\makeatletter
\newsavebox\pandoc@box
\newcommand*\pandocbounded[1]{% scales image to fit in text height/width
  \sbox\pandoc@box{#1}%
  \Gscale@div\@tempa{\textheight}{\dimexpr\ht\pandoc@box+\dp\pandoc@box\relax}%
  \Gscale@div\@tempb{\linewidth}{\wd\pandoc@box}%
  \ifdim\@tempb\p@<\@tempa\p@\let\@tempa\@tempb\fi% select the smaller of both
  \ifdim\@tempa\p@<\p@\scalebox{\@tempa}{\usebox\pandoc@box}%
  \else\usebox{\pandoc@box}%
  \fi%
}
% Set default figure placement to htbp
\def\fps@figure{htbp}
\makeatother



\ifLuaTeX
\usepackage[bidi=basic,shorthands=off,provide=*]{babel}
\else
\usepackage[bidi=default,shorthands=off,provide=*]{babel}
\fi
\ifPDFTeX
\else
\babelfont{rm}[]{DejaVu Serif}
\fi
\ifLuaTeX
  \usepackage{selnolig} % disable illegal ligatures
\fi


\setlength{\emergencystretch}{3em} % prevent overfull lines

\providecommand{\tightlist}{%
  \setlength{\itemsep}{0pt}\setlength{\parskip}{0pt}}



 

\usepackage[]{csquotes}

\IfFileExists{plex-otf.sty}{
  %% Full TeXlive
  % \usepackage[%
  %   % math,
  %   RM={Scale=0.94},SS={Scale=0.94},SScon={Scale=0.94},TT={Scale=MatchLowercase,FakeStretch=0.9},DefaultFeatures={Ligatures=Common}
  % ]{plex-otf}
}{
  %% TinyTeX
  \usepackage{libertine}
}

%%% Load theme
% https://deic.uab.cat/~iblanes/beamer_gallery/
\IfFileExists{beamerthemegotham.sty}{
  %% Full TeXlive
  \usetheme{gotham}
  \gothamset{
    numbering=totalpagenumber,
    parttocframe default=off,
    sectiontocframe default=off,
    subsectiontocframe default=off,
  }
}{
  %% TinyTeX
  \usetheme{Madrid}
}
\makeatletter
\@ifpackageloaded{caption}{}{\usepackage{caption}}
\AtBeginDocument{%
\ifdefined\contentsname
  \renewcommand*\contentsname{Содержание}
\else
  \newcommand\contentsname{Содержание}
\fi
\ifdefined\listfigurename
  \renewcommand*\listfigurename{Список иллюстраций}
\else
  \newcommand\listfigurename{Список иллюстраций}
\fi
\ifdefined\listtablename
  \renewcommand*\listtablename{Список таблиц}
\else
  \newcommand\listtablename{Список таблиц}
\fi
\ifdefined\figurename
  \renewcommand*\figurename{Рисунок}
\else
  \newcommand\figurename{Рисунок}
\fi
\ifdefined\tablename
  \renewcommand*\tablename{Таблица}
\else
  \newcommand\tablename{Таблица}
\fi
}
\@ifpackageloaded{float}{}{\usepackage{float}}
\floatstyle{ruled}
\@ifundefined{c@chapter}{\newfloat{codelisting}{h}{lop}}{\newfloat{codelisting}{h}{lop}[chapter]}
\floatname{codelisting}{Список}
\newcommand*\listoflistings{\listof{codelisting}{Листинги}}
\makeatother
\makeatletter
\makeatother
\makeatletter
\@ifpackageloaded{caption}{}{\usepackage{caption}}
\@ifpackageloaded{subcaption}{}{\usepackage{subcaption}}
\makeatother

\usepackage{bookmark}
\IfFileExists{xurl.sty}{\usepackage{xurl}}{} % add URL line breaks if available
\urlstyle{same}
\hypersetup{
  pdftitle={Выполнение лабораторной работы №6},
  pdfauthor={Карпухин Клим},
  pdflang={ru-RU},
  hidelinks,
  pdfcreator={LaTeX via pandoc}}


\title{\texorpdfstring{Выполнение лабораторной работы
№6}{Выполнение лабораторной работы №6}}
\subtitle{\texorpdfstring{Основы интерфейса взаимодействия пользователя
с системой Unix на уровне командной
строки.}{Основы интерфейса взаимодействия пользователя с системой Unix на уровне командной строки.}}
\author{Карпухин Клим}
\date{2026-03-21}

\begin{document}
\frame{\titlepage}


\section{Содержание}\label{ux441ux43eux434ux435ux440ux436ux430ux43dux438ux435}

\begin{frame}{Содержание}
\begin{enumerate}[<+->]
\tightlist
\item
  Информация о докладчике
\item
  Вводная часть и актуальность
\item
  Объект и предмет исследования
\item
  Научная новизна и практическая значимость
\item
  Цель, гипотеза и задачи
\item
  Материалы, методы и инструменты
\item
  Ход работы (этапы, скриншоты)
\item
  Результаты и анализ
\item
  Выводы
\end{enumerate}
\end{frame}

\section{Информация}\label{ux438ux43dux444ux43eux440ux43cux430ux446ux438ux44f}

\begin{frame}{Докладчик}
\protect\phantomsection\label{ux434ux43eux43aux43bux430ux434ux447ux438ux43a}
\begin{columns}[c]
\begin{column}{0.65\linewidth}
\begin{itemize}[<+->]
\tightlist
\item
  \textbf{Карпухин Клим}
\item
  Российский университет дружбы народов
\item
  Email:
  \href{mailto:1032255580@rudn.ru}{\nolinkurl{1032255580@rudn.ru}}
\item
  Роли: студент (лабораторная работа по ОС/виртуализации)
\end{itemize}
\end{column}

\begin{column}{0.35\linewidth}

\includegraphics[width=0.9\linewidth,height=\textheight,keepaspectratio]{./image/dokladchik.png}
\end{column}
\end{columns}
\end{frame}

\section{Вводная
часть}\label{ux432ux432ux43eux434ux43dux430ux44f-ux447ux430ux441ux442ux44c}

\begin{frame}[fragile]{Актуальность}
\protect\phantomsection\label{ux430ux43aux442ux443ux430ux43bux44cux43dux43eux441ux442ux44c}
\begin{itemize}[<+->]
\tightlist
\item
  Командная строка остаётся одним из основных инструментов работы в
  Linux, особенно при администрировании, разработке и настройке системы.
\item
  Умение работать с каталогами, файлами, справкой \texttt{man} и
  историей команд позволяет быстрее и точнее выполнять типовые операции
  без графического интерфейса.
\item
  Освоение базовых команд \texttt{cd}, \texttt{pwd}, \texttt{ls},
  \texttt{mkdir}, \texttt{rmdir}, \texttt{rm}, \texttt{history}
  формирует фундамент для дальнейшего изучения Unix-подобных систем.
\item
  Работа в Fedora Sway особенно хорошо демонстрирует, что даже в
  современном графическом окружении терминал остаётся незаменимым
  инструментом.
\end{itemize}
\end{frame}

\begin{frame}{Объект и предмет исследования}
\protect\phantomsection\label{ux43eux431ux44aux435ux43aux442-ux438-ux43fux440ux435ux434ux43cux435ux442-ux438ux441ux441ux43bux435ux434ux43eux432ux430ux43dux438ux44f}
\begin{itemize}[<+->]
\tightlist
\item
  \textbf{Объект:} процесс взаимодействия пользователя с операционной
  системой Linux через командную строку.
\item
  \textbf{Предмет:} базовые команды оболочки Unix/Linux, применяемые для
  навигации по файловой системе, управления каталогами, просмотра
  справки и работы с историей команд.
\end{itemize}
\end{frame}

\section{Научная новизна и практическая
значимость}\label{ux43dux430ux443ux447ux43dux430ux44f-ux43dux43eux432ux438ux437ux43dux430-ux438-ux43fux440ux430ux43aux442ux438ux447ux435ux441ux43aux430ux44f-ux437ux43dux430ux447ux438ux43cux43eux441ux442ux44c}

\begin{frame}[fragile]{Научная новизна}
\protect\phantomsection\label{ux43dux430ux443ux447ux43dux430ux44f-ux43dux43eux432ux438ux437ux43dux430}
\begin{itemize}[<+->]
\tightlist
\item
  Изучение поведения стандартных команд Unix/Linux в реальной среде
  Fedora Sway.
\item
  Сравнение разных режимов вывода команды \texttt{ls} для анализа
  структуры каталогов.
\item
  Практическая демонстрация способов работы с историей команд и
  модификации ранее введённых команд.
\end{itemize}
\end{frame}

\begin{frame}[fragile]{Практическая значимость работы}
\protect\phantomsection\label{ux43fux440ux430ux43aux442ux438ux447ux435ux441ux43aux430ux44f-ux437ux43dux430ux447ux438ux43cux43eux441ux442ux44c-ux440ux430ux431ux43eux442ux44b}
\begin{itemize}[<+->]
\tightlist
\item
  Полученные навыки позволяют уверенно работать в терминале без опоры на
  графический интерфейс.
\item
  Освоение командной строки ускоряет выполнение рутинных операций с
  файлами и каталогами.
\item
  Понимание справочной системы \texttt{man} помогает самостоятельно
  изучать новые команды и их параметры.
\item
  Навык использования \texttt{history} упрощает повторное выполнение и
  редактирование команд.
\end{itemize}
\end{frame}

\section{Цель, гипотеза и
задачи}\label{ux446ux435ux43bux44c-ux433ux438ux43fux43eux442ux435ux437ux430-ux438-ux437ux430ux434ux430ux447ux438}

\begin{frame}{Цель}
\protect\phantomsection\label{ux446ux435ux43bux44c}
Освоить базовые команды командной строки Unix/Linux, научиться работать
с файловой системой, справочной системой и историей команд.
\end{frame}

\begin{frame}[fragile]{Гипотеза}
\protect\phantomsection\label{ux433ux438ux43fux43eux442ux435ux437ux430}
Если освоить основные команды оболочки Linux и научиться пользоваться
справкой \texttt{man} и историей \texttt{history}, то работа в терминале
станет быстрее, удобнее и надёжнее.
\end{frame}

\begin{frame}[fragile]{Задачи}
\protect\phantomsection\label{ux437ux430ux434ux430ux447ux438}
\begin{itemize}[<+->]
\tightlist
\item
  Определить полный путь домашнего каталога.
\item
  Перейти в каталог \texttt{/tmp} и изучить его содержимое с помощью
  \texttt{ls}.
\item
  Проверить наличие каталога \texttt{cron} в \texttt{/var/spool}.
\item
  Создать и удалить каталоги с помощью \texttt{mkdir}, \texttt{rmdir} и
  \texttt{rm}.
\item
  Определить опции \texttt{ls} для рекурсивного просмотра и сортировки
  по времени.
\item
  Изучить справку по командам \texttt{cd}, \texttt{pwd}, \texttt{mkdir},
  \texttt{rmdir}, \texttt{rm}.
\item
  Использовать \texttt{history} для повторного выполнения и изменения
  ранее введённых команд.
\end{itemize}
\end{frame}

\section{Материалы и
методы}\label{ux43cux430ux442ux435ux440ux438ux430ux43bux44b-ux438-ux43cux435ux442ux43eux434ux44b}

\begin{frame}[fragile]{Материалы и методы}
\protect\phantomsection\label{ux43cux430ux442ux435ux440ux438ux430ux43bux44b-ux438-ux43cux435ux442ux43eux434ux44b-1}
\begin{itemize}[<+->]
\tightlist
\item
  \textbf{Операционная система:} Fedora Linux.
\item
  \textbf{Окружение рабочего стола:} Sway.
\item
  \textbf{Инструменты:}

  \begin{itemize}[<+->]
  \tightlist
  \item
    терминал;
  \item
    команда \texttt{pwd};
  \item
    команда \texttt{ls} с различными опциями;
  \item
    команда \texttt{cd};
  \item
    команды \texttt{mkdir}, \texttt{rmdir}, \texttt{rm};
  \item
    команда \texttt{man};
  \item
    команда \texttt{history}.
  \end{itemize}
\item
  \textbf{Методика:} последовательное выполнение команд в терминале с
  фиксацией результата на скриншотах и анализом вывода.
\end{itemize}
\end{frame}

\section{Ход
работы}\label{ux445ux43eux434-ux440ux430ux431ux43eux442ux44b}

\section{Ход работы --- выполнение
заданий}\label{ux445ux43eux434-ux440ux430ux431ux43eux442ux44b-ux432ux44bux43fux43eux43bux43dux435ux43dux438ux435-ux437ux430ux434ux430ux43dux438ux439}

\begin{frame}{Этап 1: Определение домашнего каталога}
\protect\phantomsection\label{ux44dux442ux430ux43f-1-ux43eux43fux440ux435ux434ux435ux43bux435ux43dux438ux435-ux434ux43eux43cux430ux448ux43dux435ux433ux43e-ux43aux430ux442ux430ux43bux43eux433ux430}
Определил полное имя моего домашнего каталога. (рис.~\ref{fig-001}).

\begin{figure}

\centering{

\includegraphics[width=0.7\linewidth,height=\textheight,keepaspectratio]{image/1.png}

}

\caption{\label{fig-001}Определение полного имени моего домашнего
каталога}

\end{figure}%
\end{frame}

\begin{frame}[fragile]{Этап 2: Работа с каталогом \texttt{/tmp}}
\protect\phantomsection\label{ux44dux442ux430ux43f-2-ux440ux430ux431ux43eux442ux430-ux441-ux43aux430ux442ux430ux43bux43eux433ux43eux43c-tmp}
\begin{block}{2.1 Переход в каталог \texttt{/tmp}}
\protect\phantomsection\label{ux43fux435ux440ux435ux445ux43eux434-ux432-ux43aux430ux442ux430ux43bux43eux433-tmp}
Перешёл в каталог \texttt{/tmp}. (рис.~\ref{fig-002}).

\begin{figure}

\centering{

\includegraphics[width=0.7\linewidth,height=\textheight,keepaspectratio]{image/2.png}

}

\caption{\label{fig-002}Переход в каталог /tmp}

\end{figure}%
\end{block}

\begin{block}{2.2 Просмотр содержимого каталога \texttt{/tmp}}
\protect\phantomsection\label{ux43fux440ux43eux441ux43cux43eux442ux440-ux441ux43eux434ux435ux440ux436ux438ux43cux43eux433ux43e-ux43aux430ux442ux430ux43bux43eux433ux430-tmp}
Вывел на экран содержимое каталога \texttt{/tmp} командой \texttt{ls}.
(рис.~\ref{fig-003}).

\begin{figure}

\centering{

\includegraphics[width=0.7\linewidth,height=\textheight,keepaspectratio]{image/3.png}

}

\caption{\label{fig-003}Вывод содержимого каталога /tmp командой
\texttt{ls}}

\end{figure}%

Также были использованы различные опции команды \texttt{ls} для
получения более подробной информации о содержимом каталога
\texttt{/tmp}:

\begin{verbatim}
ls -a
ls -l
ls -F
ls -alF
\end{verbatim}

Обычная команда \texttt{ls} показывает только имена файлов и каталогов.
Опция \texttt{-a} добавляет скрытые файлы и каталоги, имена которых
начинаются с точки. Опция \texttt{-l} выводит подробную информацию:
права доступа, число ссылок, владельца, группу, размер и дату изменения.
Опция \texttt{-F} добавляет символ, указывающий тип объекта: \texttt{/}
для каталога, \texttt{*} для исполняемого файла, \texttt{@} для
символической ссылки. Комбинированный вариант \texttt{ls\ -alF}
объединяет все эти возможности.
\end{block}

\begin{block}{2.3 Проверка наличия каталога \texttt{cron}}
\protect\phantomsection\label{ux43fux440ux43eux432ux435ux440ux43aux430-ux43dux430ux43bux438ux447ux438ux44f-ux43aux430ux442ux430ux43bux43eux433ux430-cron}
Определил, есть ли в каталоге \texttt{/var/spool} подкаталог с именем
\texttt{cron}. (рис.~\ref{fig-004}).

\begin{figure}

\centering{

\includegraphics[width=0.7\linewidth,height=\textheight,keepaspectratio]{image/4.png}

}

\caption{\label{fig-004}Определение наличия каталога /var/spool/cron}

\end{figure}%
\end{block}

\begin{block}{2.4 Просмотр домашнего каталога}
\protect\phantomsection\label{ux43fux440ux43eux441ux43cux43eux442ux440-ux434ux43eux43cux430ux448ux43dux435ux433ux43e-ux43aux430ux442ux430ux43bux43eux433ux430}
Перешёл в мой домашний каталог и вывел на экран его содержимое.
(рис.~\ref{fig-005}).

\begin{figure}

\centering{

\includegraphics[width=0.7\linewidth,height=\textheight,keepaspectratio]{image/5.png}

}

\caption{\label{fig-005}Переход в домашний каталог, вывод содержимого}

\end{figure}%

Вывод \texttt{ls\ -l} позволил определить владельца файлов и
подкаталогов. В соответствующем столбце отображается имя пользователя,
которому принадлежат объекты файловой системы.
\end{block}
\end{frame}

\begin{frame}[fragile]{Этап 3: Создание и удаление каталогов}
\protect\phantomsection\label{ux44dux442ux430ux43f-3-ux441ux43eux437ux434ux430ux43dux438ux435-ux438-ux443ux434ux430ux43bux435ux43dux438ux435-ux43aux430ux442ux430ux43bux43eux433ux43eux432}
\begin{block}{3.1 Создание каталога \texttt{newdir}}
\protect\phantomsection\label{ux441ux43eux437ux434ux430ux43dux438ux435-ux43aux430ux442ux430ux43bux43eux433ux430-newdir}
В домашнем каталоге создал новый каталог с именем \texttt{newdir}.
(рис.~\ref{fig-006}).

\begin{figure}

\centering{

\includegraphics[width=0.7\linewidth,height=\textheight,keepaspectratio]{image/6.png}

}

\caption{\label{fig-006}Создание каталога newdir}

\end{figure}%
\end{block}

\begin{block}{3.2 Создание каталога \texttt{morefun}}
\protect\phantomsection\label{ux441ux43eux437ux434ux430ux43dux438ux435-ux43aux430ux442ux430ux43bux43eux433ux430-morefun}
В каталоге \texttt{\textasciitilde{}/newdir} создал новый каталог с
именем \texttt{morefun}. (рис.~\ref{fig-007}).

\begin{figure}

\centering{

\includegraphics[width=0.7\linewidth,height=\textheight,keepaspectratio]{image/7.png}

}

\caption{\label{fig-007}Создание каталога morefun}

\end{figure}%
\end{block}

\begin{block}{3.3 Создание трёх каталогов одной командой}
\protect\phantomsection\label{ux441ux43eux437ux434ux430ux43dux438ux435-ux442ux440ux451ux445-ux43aux430ux442ux430ux43bux43eux433ux43eux432-ux43eux434ux43dux43eux439-ux43aux43eux43cux430ux43dux434ux43eux439}
В домашнем каталоге создал одной командой три новых каталога с именами
\texttt{letters}, \texttt{memos}, \texttt{misk}. (рис.~\ref{fig-008}).

\begin{figure}

\centering{

\includegraphics[width=0.7\linewidth,height=\textheight,keepaspectratio]{image/8.png}

}

\caption{\label{fig-008}Создание трёх каталогов одной командой}

\end{figure}%

Затем удалил эти каталоги одной командой. (рис.~\ref{fig-009}).

\begin{figure}

\centering{

\includegraphics[width=0.7\linewidth,height=\textheight,keepaspectratio]{image/9.png}

}

\caption{\label{fig-009}Удаление трёх каталогов одной командой}

\end{figure}%
\end{block}

\begin{block}{3.4 Попытка удалить каталог \texttt{newdir} командой
\texttt{rm}}
\protect\phantomsection\label{ux43fux43eux43fux44bux442ux43aux430-ux443ux434ux430ux43bux438ux442ux44c-ux43aux430ux442ux430ux43bux43eux433-newdir-ux43aux43eux43cux430ux43dux434ux43eux439-rm}
Попробовал удалить ранее созданный каталог
\texttt{\textasciitilde{}/newdir} командой \texttt{rm}. Каталог не был
удалён. (рис.~\ref{fig-010}).

\begin{figure}

\centering{

\includegraphics[width=0.7\linewidth,height=\textheight,keepaspectratio]{image/10.png}

}

\caption{\label{fig-010}Попытка удалить каталог newdir командой rm}

\end{figure}%
\end{block}

\begin{block}{3.5 Удаление каталога \texttt{morefun}}
\protect\phantomsection\label{ux443ux434ux430ux43bux435ux43dux438ux435-ux43aux430ux442ux430ux43bux43eux433ux430-morefun}
Удалил каталог \texttt{\textasciitilde{}/newdir/morefun} из домашнего
каталога. Каталог был удалён. (рис.~\ref{fig-011}).

\begin{figure}

\centering{

\includegraphics[width=0.7\linewidth,height=\textheight,keepaspectratio]{image/11.png}

}

\caption{\label{fig-011}Удаление каталога morefun}

\end{figure}%
\end{block}
\end{frame}

\begin{frame}[fragile]{Этап 4: Работа со справкой \texttt{man}}
\protect\phantomsection\label{ux44dux442ux430ux43f-4-ux440ux430ux431ux43eux442ux430-ux441ux43e-ux441ux43fux440ux430ux432ux43aux43eux439-man}
\begin{block}{4.1 Опция для рекурсивного просмотра}
\protect\phantomsection\label{ux43eux43fux446ux438ux44f-ux434ux43bux44f-ux440ux435ux43aux443ux440ux441ux438ux432ux43dux43eux433ux43e-ux43fux440ux43eux441ux43cux43eux442ux440ux430}
С помощью команды \texttt{man\ ls} нашёл опцию для просмотра содержимого
не только указанного каталога, но и подкаталогов, входящих в него:
\texttt{ls\ -R}. (рис.~\ref{fig-012}).

\begin{figure}

\centering{

\includegraphics[width=0.7\linewidth,height=\textheight,keepaspectratio]{image/12.png}

}

\caption{\label{fig-012}Опция для рекурсивного просмотра содержимого
каталогов}

\end{figure}%
\end{block}

\begin{block}{4.2 Опции для сортировки по времени}
\protect\phantomsection\label{ux43eux43fux446ux438ux438-ux434ux43bux44f-ux441ux43eux440ux442ux438ux440ux43eux432ux43aux438-ux43fux43e-ux432ux440ux435ux43cux435ux43dux438}
С помощью команды \texttt{man\ ls} нашёл опции, позволяющие
отсортировать по времени последнего изменения выводимый список
содержимого каталога с развёрнутым описанием файлов: \texttt{ls\ -lt}.
(рис.~\ref{fig-013}).

\begin{figure}

\centering{

\includegraphics[width=0.7\linewidth,height=\textheight,keepaspectratio]{image/13.png}

}

\caption{\label{fig-013}Опции для сортировки по времени}

\end{figure}%
\end{block}

\begin{block}{4.3 Просмотр описания основных команд}
\protect\phantomsection\label{ux43fux440ux43eux441ux43cux43eux442ux440-ux43eux43fux438ux441ux430ux43dux438ux44f-ux43eux441ux43dux43eux432ux43dux44bux445-ux43aux43eux43cux430ux43dux434}
Использовал команду \texttt{man} для просмотра описания следующих
команд: \texttt{cd}, \texttt{pwd}, \texttt{mkdir}, \texttt{rmdir},
\texttt{rm}. (рис.~\ref{fig-014}).

\begin{figure}

\centering{

\includegraphics[width=0.7\linewidth,height=\textheight,keepaspectratio]{image/14.png}

}

\caption{\label{fig-014}Просмотр описания команды cd}

\end{figure}%

\begin{itemize}[<+->]
\tightlist
\item
  \texttt{cd} --- переход в другой каталог. Без аргументов переходит в
  домашний каталог, \texttt{cd\ ..} поднимает на уровень выше.
\item
  \texttt{pwd} --- выводит абсолютный путь текущего каталога.
\item
  \texttt{mkdir} --- создаёт каталоги. Опция \texttt{-p} позволяет
  создавать промежуточные каталоги.
\item
  \texttt{rmdir} --- удаляет только пустые каталоги.
\item
  \texttt{rm} --- удаляет файлы. Для удаления каталогов используется
  \texttt{-r}, а для подтверждения удаления --- \texttt{-i}.
\end{itemize}
\end{block}
\end{frame}

\begin{frame}[fragile]{Этап 5: Использование истории команд}
\protect\phantomsection\label{ux44dux442ux430ux43f-5-ux438ux441ux43fux43eux43bux44cux437ux43eux432ux430ux43dux438ux435-ux438ux441ux442ux43eux440ux438ux438-ux43aux43eux43cux430ux43dux434}
Используя информацию, полученную при помощи команды \texttt{history},
выполнил модификацию и исполнение нескольких команд из буфера команд.
(рис.~\ref{fig-015}).

\begin{figure}

\centering{

\includegraphics[width=0.7\linewidth,height=\textheight,keepaspectratio]{image/15.png}

}

\caption{\label{fig-015}Использование команды history для повторного
выполнения команд}

\end{figure}%
\end{frame}

\section{Результаты и
анализ}\label{ux440ux435ux437ux443ux43bux44cux442ux430ux442ux44b-ux438-ux430ux43dux430ux43bux438ux437}

\begin{frame}[fragile]{Анализ достигнутых результатов}
\protect\phantomsection\label{ux430ux43dux430ux43bux438ux437-ux434ux43eux441ux442ux438ux433ux43dux443ux442ux44bux445-ux440ux435ux437ux443ux43bux44cux442ux430ux442ux43eux432}
\begin{itemize}[<+->]
\tightlist
\item
  Определён полный путь домашнего каталога.
\item
  Освоен переход между каталогами и просмотр содержимого файловой
  системы.
\item
  Изучены различия между вариантами команды \texttt{ls}.
\item
  Проверено наличие системного каталога \texttt{cron} в
  \texttt{/var/spool}.
\item
  Получен практический опыт создания и удаления каталогов.
\item
  Использована справка \texttt{man} для поиска нужных опций команд.
\item
  Опробована работа с историей команд \texttt{history}.
\end{itemize}
\end{frame}

\begin{frame}[fragile]{Практическая значимость результатов}
\protect\phantomsection\label{ux43fux440ux430ux43aux442ux438ux447ux435ux441ux43aux430ux44f-ux437ux43dux430ux447ux438ux43cux43eux441ux442ux44c-ux440ux435ux437ux443ux43bux44cux442ux430ux442ux43eux432}
\begin{itemize}[<+->]
\tightlist
\item
  Я научился уверенно ориентироваться в файловой системе Linux.
\item
  Освоил базовые операции, необходимые для повседневной работы в
  терминале.
\item
  Получил представление о том, как быстро искать информацию о командах и
  их параметрах.
\item
  На практике увидел, чем отличаются команды \texttt{rm} и
  \texttt{rmdir}, а также как удобно использовать историю команд.
\end{itemize}
\end{frame}

\section{Выводы}\label{ux432ux44bux432ux43eux434ux44b}

\begin{frame}{Общее заключение}
\protect\phantomsection\label{ux43eux431ux449ux435ux435-ux437ux430ux43aux43bux44eux447ux435ux43dux438ux435}
\begin{itemize}[<+->]
\tightlist
\item
  В ходе лабораторной работы я освоил основные приёмы работы в командной
  строке Unix/Linux.
\item
  Я научился перемещаться по каталогам, просматривать содержимое
  каталогов в разных режимах, создавать и удалять каталоги, пользоваться
  справочной системой и историей команд.
\item
  Полученные навыки являются базой для дальнейшей работы в Linux и
  позволяют эффективнее выполнять задачи в терминале.
\end{itemize}
\end{frame}

\begin{frame}[fragile]{Выводы}
\protect\phantomsection\label{ux432ux44bux432ux43eux434ux44b-1}
\begin{enumerate}[<+->]
\tightlist
\item
  Командная строка является основным инструментом управления системой
  Unix/Linux.
\item
  Команда \texttt{pwd} позволяет определить абсолютный путь текущего
  каталога.
\item
  Команда \texttt{ls} с различными опциями даёт удобный способ анализа
  содержимого каталогов.
\item
  Команды \texttt{mkdir}, \texttt{rmdir} и \texttt{rm} позволяют
  создавать и удалять каталоги при соблюдении их особенностей.
\item
  Команда \texttt{man} помогает самостоятельно изучать синтаксис и опции
  команд.
\item
  Команда \texttt{history} ускоряет работу за счёт повторного
  использования и изменения ранее введённых команд.
\end{enumerate}
\end{frame}




\end{document}
